\begin{center}
    {\LARGE \textbf{2013分析化学}} \\[2ex]
    {\large 时量180分钟 \quad 满分90分}
\end{center}

\vspace{0.5cm}
\noindent\textbf{注意事项:}
\begin{enumerate}[label=\arabic*., parsep=0pt, itemsep=0pt]
    \item 答题前,考生先将自己的姓名、学号、考场号、座位号填写在答题卡上,将条形码准确粘贴在考生信息条形码粘贴区。
    \item 考试结束后,将本试卷与答题纸一并收回。
    \item 回答各个问题时,将答案写在答题纸上,写在本试卷上无效。
\end{enumerate}

\vspace{0.5cm}
\noindent\textbf{一、简答题(35 pts.)}

\begin{enumerate}[label=\arabic*., itemsep=1.5ex]
    \item 要使置信度$90\%$的平均值的置信区间不超过$\pm s$和$\pm 2s$，问至少各应平行测定几次？结果说明什么？
    \\
    已知：
    \begin{center}
    \begin{tabular}{l|ccccc}
    $f$ & $2$ & $3$ & $4$ & $5$ & $6$ \\
    \hline
    $t_{0.10}$ & $6.31$ & $2.92$ & $2.35$ & $2.13$ & $2.02$ \\
    \end{tabular}
    \end{center}
    \item 为什么不能用直接法配制$\ce{KMnO4}$标准溶液？配制和保存$\ce{KMnO4}$标准溶液时应注意些什么问题？
    \item 重量分析中为什么必须加入适当过量的沉淀剂？为什么又不能过量太多。
    \item 简述晶型沉淀的实验条件，并解释采取这些措施的理由。
    \item 计算在$25^\circ\mathrm{C}$ $[\ce{I^-}]=1\ \mathrm{mol/L}$时，$\ce{Cu^{2+}/Cu^+}$电对的条件电位(忽略离子强度的影响)
    
    [已知：$E^\ominus_{\ce{Cu^{2+}/Cu^+}}=0.159\ \mathrm{V}$，$K_{\mathrm{sp}}(\ce{CuI})=1.1\times10^{-12}$]
    \item 用$0.1224\ \mathrm{mol/L}\ \ce{HCl}$滴定可能含有$\ce{NaOH}$,$\ce{Na2CO3}$,$\ce{NaHCO3}$的溶液，滴至酚酞终点消耗$34.66\ \mathrm{mL}$，滴至溴甲酚绿终点消耗$41.24\ \mathrm{mL}$，判断溶液的组成，并计算各组分的质量。已知$M_r(\ce{NaOH})=40$，$M_r(\ce{Na2CO3})=106$，$M_r(\ce{NaHCO3})=84$
    \item 若某物质在两相中的分配比为$17$，今有$50\ \mathrm{mL}$的水溶液，应用$10\ \mathrm{mL}$的萃取剂连续萃取几次，才能使总萃取率达$95\%$以上？
\end{enumerate}

\vspace{0.5cm}
\noindent\textbf{二、计算题(55 pts. total)}

\begin{enumerate}[label=\arabic*., itemsep=1.5ex]
    \item (10 pts.) 在$\mathrm{pH}=9.0$的$\ce{NH3-NH4Cl}$缓冲液中$[\ce{NH3}]=0.10\ \mathrm{mol/L}$通入$\ce{H2S}$至饱和（浓度约为$0.010\ \mathrm{mol/L}$），求$\ce{Ag2S}$在此溶液中的溶解度，
    
    [$\ce{H2S}$ $K_{a1}=1.3\times10^{-7}$，$K_{a2}=7.1\times10^{-15}$；$\ce{Ag^+} \sim \ce{NH3}$络合物$\lg\beta_1$ $\lg\beta_2$分别为$3.4$和$7.4$，$K_{\mathrm{sp}}(\ce{Ag2S})=2\times10^{-50}$]
    \item (10 pts.) 某磷酸盐缓冲溶液，每升含$0.080\ \mathrm{mol}\ \ce{Na2HPO4}$和$0.020\ \mathrm{mol}\ \ce{Na3PO4}$。今有$1.0\ \mathrm{mmol}$有机化合物$\ce{RNHOH}$在$100\ \mathrm{mL}$上述磷酸盐缓冲溶液中进行电解氧化，其反应如下
    \[
    \ce{RNHOH + H2O -> RNO2 + 4H+ + 4e^-}
    \]
    当氧化反应进行完全后，计算缓冲溶液的$\mathrm{pH}$。［已知$\ce{H3PO4}$的$\mathrm{p}K_{a1} \sim \mathrm{p}K_{a3}$分别为$2.12$，$7.20$，$12.36$］
    \item (12 pts.) 用$0.2\ \mathrm{mol/L}\ \ce{NaOH}$滴定浓度均为$0.2\ \mathrm{mol/L}$的$\ce{HCl}$和$\ce{NH4Cl}$混合液中的$\ce{HCl}$，计算滴定突跃，用$\mathrm{MR}$为指示剂，当终点$\mathrm{pH}_{ep}=6.2$时求滴定误差。
    
    ［已知$\mathrm{p}K_b(\ce{NH3})=4.74$］
    \item (12 pts.) 在$\mathrm{pH}=5.0$ $[\ce{HAc}]=0.2\ \mathrm{mol/L}$，$[\ce{Ac^-}]=0.4\ \mathrm{mol/L}$的介质条件下，用$2\times10^{-3}\ \mathrm{mol/L}\ \mathrm{EDTA}$滴定同浓度$\ce{Pb^{2+}}$，若终点$\mathrm{pPb_{ep}}=7.0$，求滴定误差。
    
    ［已知$\lg K_{\ce{PbY}}=18.0$，$\mathrm{pH}=5.0$时$\lg\alpha_{Y(\mathrm{H})}=6.4$，$\ce{Pb-Ac}$的$\lg\beta_1=1.9$，$\lg\beta_2=3.3$］
    \item (11 pts.) 称取某合金钢试样$0.2000\ \mathrm{g}$酸溶后其中的钒被氧化为$\ce{VO_2^+}$并使$\ce{VO_2^+}$与钽试剂反应生成有色螯合物，定容为$100\ \mathrm{mL}$，然后取出部分溶液用等体积的$\ce{CHCl3}$萃取一次（设分配比$D=10$）有机相在$530\ \mathrm{nm}$处有最大吸收，$\varepsilon_{530}=5.7\times10^4\ \mathrm{L/(mol\cdot cm)}$，若使用$1\ \mathrm{cm}$的比色皿，测得吸光度$A=0.570$，计算试样中$\mathrm{V\%}$（已知$A_r(\mathrm{V})=50.94$）
\end{enumerate}